% Ad M. Brutum 1.16 (ad-brutum-01-16) --- English translation
% Translated by Claude (Anthropic) from the Latin (Perseus).
% Style guide: STYLE.md.

\ciceroLetterOpener{Brutus Ciceroni salutem}

\ciceroSection{1} I have read a fragment of that little letter of yours to Octavius, the one Atticus sent on to me. Your concern and your care for my safety gave me no new pleasure: it is not only usual but a daily thing to hear of something you have said or done loyally and with honour for our standing. But that same part of the letter, the part written to Octavius about us, struck me with as much grief as my mind is able to hold. For you thank him for the state in terms so suppliant, so abject --- what am I to write? I am ashamed of our condition and our fortune, but it must be written all the same: you commend our safety to him. What death could be more ruinous than such safety? You make it plain, in so many words, that the despotism has not been abolished, only the master changed. Look back over your own words, and dare to deny that these are the prayers of a slave addressed to a king. There is one thing, you say, that is asked and looked for from him: that he should be willing to keep safe those citizens of whom good men and the Roman people think well. And what if he is not willing? Shall we then cease to exist? Yet it is better not to exist than to exist by his leave.

\ciceroSection{2} I myself, upon my word, do not believe that all the gods are so set against the safety of the Roman people that Octavius must be petitioned for the safety of any citizen at all --- I will not say for the liberators of the whole world. It is a pleasure to speak in grand terms, and it is surely fitting to do so before men who do not know what is to be feared for each man and what is to be sought from each. Do you, Cicero, admit that Octavius has this power, and are you his friend? Or, if you hold me dear, do you wish to see me at Rome, when I had to be commended to that boy before I could so much as be there? Why do you thank him, if you think he must be entreated to be willing that we be safe, and to suffer it? Or is this to be reckoned a kindness, that he preferred to be the one of whom such things must be begged, rather than Antony? To the avenger of another's despotism --- not to its inheritor --- does anyone make supplication that men who have deserved supremely well of the state be allowed to live?

\ciceroSection{3} It is this very weakness and despair --- the fault of which lies no more with you than with everyone else --- that drove Caesar into his lust for kingship, that persuaded Antony, after Caesar's death, to try to seize the place of the man who had been killed, and that now has lifted up this boy of yours, so that you should judge that the safety of such men must be won by prayers, and that we shall scarcely even now be safe through the mercy of one man --- a man, barely --- and by no other thing whatever. But if we had remembered that we are Romans, the basest of men would no more boldly long to be our masters than we would prevent it; nor would Antony have been so much inflamed by the desire for Caesar's throne as he was deterred from it by Caesar's death.

\ciceroSection{4} You yourself --- a consular, the avenger of such great crimes, crimes which, now that they are crushed, I fear have been put off by you only for a short time --- how can you look upon what you have achieved, and at the same time either approve these present doings or endure them so abjectly and so easily that you wear the look of one who approves? And what private quarrel have you with Antony? Surely it was this: that he demanded these very things --- that safety be sought from him, that we should hold our security as a thing granted on sufferance from the very men from whom he himself had received his freedom, that the disposal of the state should be his to make. Was it for this that you thought arms must be sought, by which he might be prevented from being our master --- so that, when he was prevented, we might beg the same of another, who would let himself be set up in his place? Or was it that the state might be its own, free and master of itself? Unless, perhaps, what we refused was not slavery but the terms of our slavery. And yet under a good master, Antony, we could not only have borne our fortune, but might even have enjoyed kindnesses and honours, sharing in them as fully as we pleased. For what would he refuse to men whose compliance he saw to be the greatest safeguard of his own despotism? But nothing was worth so much that we should sell our honour and our liberty for it.

\ciceroSection{5} This very boy, whom Caesar's name seems to spur on against Caesar's killers --- how much would he give, if there were room for a bargain, to have as much power as he assuredly will have, with us to back him, since we are willing to live, to keep our money, and to be called consulars! But will it have been for nothing that the man perished at whose death we rejoiced --- if, though he was dead, we were going to be slaves no less than before? Is no thought given to that? But may all the gods and goddesses sooner snatch everything from me than that resolve of mine, by which I would not grant even to the heir of the man I killed --- much less to my own father, were he to come back to life --- what I would not endure in the man himself: that, with my consent, he should have more power than the laws and the Senate. Or are you persuaded that the rest of us will be free under a man without whose consent there is no place for us in this state of ours? And how, besides, can what you ask be brought about, that you should obtain it? For you ask that he be willing that we be safe. Do we then seem to you about to receive safety, when we shall have received life? And how can we receive that life, unless we first lay down our dignity and our liberty?

\ciceroSection{6} Or do you think that to live at Rome is to be safe? It is the thing, not the place, that must guarantee that to me. I was not safe while Caesar lived --- not until I had carried out that deed; and nowhere can I be an exile, so long as I hate slavery and the endurance of insults worse than all other evils. Is this not to have fallen back into the same darkness: if from the man who has taken the name of tyrant to himself --- when in the cities of the Greeks the children of tyrants, once the tyrants are overthrown, are visited with the very same punishment --- it is begged that the avengers and overthrowers of despotism be kept safe? Could I wish to see such a state, or think any state worth the name, that cannot take back its liberty even when it is handed over and pressed upon it --- that fears in a boy the name of the king who has been removed more than it trusts in itself, when it sees that very king, who held the greatest power, removed by the courage of a few? As for me, never again commend me to your Caesar --- nor even yourself, if you will hear me. You set a very high price on the few years that your age admits of, if for their sake you are going to play the suppliant to that boy.

\ciceroSection{7} And then, the very thing you did, and are still doing, so admirably in the case of Antony --- take care that it does not turn from a credit to your greatness of spirit into a reputation for cowardice. For if Octavius pleases you, the man from whom our safety is to be sought, you will seem not to have fled a master, but to have looked for a more agreeable one. That you praise him for what he has done so far, I wholly approve; for those things are praiseworthy, provided he undertook those actions against another's power and not on behalf of his own. But when you judge that so much not only may be allowed him, but must even be conferred on him by you yourself, that he must be begged not to refuse to let us be safe, you set the price too high (for you are bestowing on him the very thing the state was thought to hold through him); nor does it cross your mind that, if Octavius deserves any honours at all because he is waging war against Antony, the Roman people will never confer on those who have cut out the evil of which these are the remnants anything that could equal their desert, even if it heaped up all rewards at once.

\ciceroSection{8} And see how much more diligently men fear than they remember. Because Antony is alive and under arms, while in Caesar's case what could and ought to have been done is finished and cannot now be restored to its first state --- Octavius is the man whose judgement upon us the Roman people is to await; we are the men for whose safety one single man, it seems, must be entreated. I, for my part --- to come back to that point --- am the man who would not only refuse to play the suppliant, but would even hold in check those who demand to be supplicated. Either I shall keep far away from slaves, and shall judge that Rome is for me wherever it shall be permitted to be free; or I shall pity you, you for whom neither age nor honours nor another man's courage has been able to lessen the sweetness of mere living.

\ciceroSection{9} For my part, I shall think myself happy enough, if only this resolve holds firm and unbroken, so that I may count the gratitude owed to my devotion as paid in full. For what is better than to be content with the memory of right deeds and with liberty, and to think nothing of human things? But I shall surely not give way to those who give way, nor be conquered by those who wish to be conquered; I shall make trial of everything and attempt everything, and I shall not cease to drag our state away from servitude. If the fortune that ought to follow does follow, we shall all rejoice; if not, I shall rejoice all the same. For with what deeds or thoughts could this life of mine better be spent than with those that have been directed to the freeing of my fellow citizens?

\ciceroSection{10} You, Cicero, I beg and exhort: do not grow weary, do not despair, but always, in fending off present evils, look ahead to those to come as well, lest they steal in unless they are met beforehand. That brave and free spirit, with which both as consul and now as consular you have championed the state --- believe that without steadiness and evenness it is nothing. For I confess that the condition of proven virtue is harder than that of virtue untried. Good deeds we exact as our due; and when they turn out otherwise, we reproach the men who have disappointed us with a hostile heart. And so, that Cicero resists Antony, though it is worthy of the highest praise, still wins no one's wonder, since that consul is rightly seen to live up to this consular.

\ciceroSection{11} But let this same Cicero turn against others the judgement that he has aimed, with such firmness and grandeur, at driving out Antony --- and he will not only have robbed himself of the glory of the time to come, but will compel even what is past to fade away. For nothing is great in itself unless the reasoned judgement behind it stands plain. And indeed no one is more fitted to love the state and to be the defender of its liberty, whether by his genius, or by his deeds, or by the wish and the urgent demand of all men. Wherefore Octavius is not to be begged to be willing that we be safe; rather, rouse yourself, and believe that the state in which you have done your greatest deeds will be free and honourable, if only the people have leaders to resist the designs of wicked men.
